\documentclass[11pt]{article}

\newcommand{\cnum}{Math 245B}
\newcommand{\ced}{Winter 2026}
\newcommand{\ctitle}[4]{\title{\vspace{-0.5in}\cnum, \ced\\Assignment #1: #2}\author{\vspace{-0.35in}\\#3, #4}}
\usepackage{enumitem}
\usepackage[usenames,dvipsnames,svgnames,table,hyperref]{xcolor}
\usepackage{amsmath, amsfonts}
\usepackage{graphicx} % Required for inserting images
\usepackage{float}
\usepackage{listings}
\usepackage{xcolor}
\usepackage{hyperref}
\usepackage{comment}

\renewcommand*{\theenumi}{\alph{enumi}}
\renewcommand*\labelenumi{(\theenumi)}
\renewcommand*{\theenumii}{\roman{enumii}}
\renewcommand*\labelenumii{\theenumii.}

\definecolor{codegreen}{rgb}{0,0.6,0}
\definecolor{codegray}{rgb}{0.5,0.5,0.5}
\definecolor{codepurple}{rgb}{0.58,0,0.82}
\definecolor{backcolour}{rgb}{0.95,0.95,0.92}
\definecolor{header}{HTML}{205459}
\definecolor{solution}{HTML}{204d52}

\newcommand{\solution}[1]{{{\textcolor{header}{Solution:} \textcolor{solution}{#1}}}}
\setlist[enumerate]{label=(\alph*)}

\lstdefinestyle{mystyle}{
    backgroundcolor=\color{backcolour},
    commentstyle=\color{codegreen},
    keywordstyle=\color{magenta},
    numberstyle=\tiny\color{codegray},
    stringstyle=\color{codepurple},
    basicstyle=\ttfamily\footnotesize,
    breakatwhitespace=false,
    breaklines=true,
    captionpos=b,
    keepspaces=true,
    numbers=left,
    numbersep=5pt,
    showspaces=false,
    showstringspaces=false,
    showtabs=false,
    tabsize=2
}


\begin{document}
\ctitle{\#1}{}{Asher Christian}{006-150-286}
\date{}
\maketitle
\[
A \in \sigma(X_1,\dots,X_n)
\] 
\[
\tilde A = \pi \circ A, \qquad \pi(1,\dots,n,n+1,\dots,2n) = (n,n+1,\dots,2n,1,2,\dots,n)
\] 
$E \in \mathcal{E}$
\[
\mathbb{P}(\vec X \in E \triangle A ) = \mathbb{P}(\vec X \in E \triangle \tilde A)
\] 
For example $X_1,X_2$ iid
\[
    \mathbb{P}(X_1 > X_2) = \tilde {\mathbb{P}}(X_2 > X_1) =^{\text{i.i.d}} \mathbb{P}(X_2 > X_1)
\] 
\[
    \mathbb{P}(\vec X \in E \triangle A) = \tilde{ \mathbb{P}}\left( \pi \circ \vec X \in E \triangle A\right) = \tilde{ \mathbb{P}} (\vec X \in E \triangle \tilde A) =^{\text{i.i.d}} \mathbb{P}(\vec X \in E \triangle \tilde A)
\] 
Let $ \mathbb{P}(\vec X \in E \triangle A) = \mathbb{P}(\vec X \in E \triangle \tilde A) = \delta$ but $A \perp \tilde A$ implies
 \[
     \mathbb{P}(E) = o(\delta) \;\; \text{or} \;\; 1 - o(\delta)
\] 
for all $\delta$ so we get
\[
    \mathbb{P}(E) = 0 \;\; \text{or} \;\; 1
\] 
impleis the $H-S$  $0-1$ Law

\section{Week 2}
Tightness + convergence in Probability implies convergence in $L_1$. For all $\epsilon > 0$ there
exists $M$ such that
 \[
     \sup_{n} \mathbb{P}(|X_n| > M) = \sup_{n} \mathbb{E}(1_{|X_n|} > M) \le \epsilon
\] \\
Uniformly integrable , $\forall \epsilon > 0$ there exists $M$ such that 
 \[
     \sup_{n} \mathbb{E}(|X_n| 1_{|X_n| > M}) \le \epsilon
\] 

\section{Theorem}
\emph{
    $\{X_n\} : \Omega \rightarrow (S, \mathcal{S})$ $(S, \mathcal{S})$ is a Polish space.
    Tight:  $\forall \epsilon > 0 \;\; \exists \; \text{compact} \; K \rightarrow \sup_{n} \mathbb{P}(X_n \in K) \le \epsilon$.
    Tight $\iff$ relatively compact wrt Weak topologoy or equivalently  $\exists \; X, n_k$ such that  $X_{n_k} \rightarrow X $ in distribution.
}\\
Proof when $ S = \mathbb{R}$: \\
Tight $\implies$ subsequence\\
$\{X_n\}$ random variables 
$F_{X_n}(a) = \mathbb{P}(X_n \le a)$
Recall $X_n \rightarrow X$ in distribution iff $\forall a$ if  $f_X(a)$ is a cont point then  $F_{X_n}(a) \rightarrow F_X(a)$.
$ \mathbb{Q} \subset \mathbb{R}$. $q_1,q_2,\dots$ $n_1,n_2,\dots$ such that
$F_{X_{n_1(k)}}(q_1) \rightarrow a_1$ and $F_{X_{n_l(k)}}(q_l) \rightarrow a_l$ and $n_l(k) \subset n_{l-1}(k)$ subsequence of the previous subsequence
then $F_{X_{n_k(k)}}$ converges for each $q_i$ simplify and rewrite as  $n_k$. Let  $f : \mathbb{Q} \rightarrow [0,1]$ be such that $f(q_i) = a_i$ we then extend this function to a function on all of $ \mathbb{R}$ and we use tightness
to

\section{Wednesday}
Continuing on tightness $\iff$ relatively compact on weak topology.
For any subsequene, exists  sub-sub sequence converging to $X$ in distribution. We proved the forward direction yesterday so, today we will prove the
other direction. If not tight then there exists $\epsilon$ and $X_{n_k}$ such that  $ \mathbb{P}(|X_{n_k}| > k) \ge \epsilon$. If $X_{n_{k_m}} \rightarrow^{d} X$ in distribution contradiction in tail.
$\exists M $ such that  $ \mathbb{P}(|X| \ge M) \le \frac{\epsilon}{2}$
\[
    \liminf_{  } \mathbb{P}(|X_{n_{k_m}}| \ge M) \ge \epsilon
\] 

\section{Lemma}
\emph{
$\{X_n\}$ If  $\exists \phi: \mathbb{R}^{+} \rightarrow \mathbb{R}$ 
\[
    \lim_{a\to +\infty}\phi(a) = +\infty \qquad \sup_{n}\mathbb{E} \phi(|X_n|) < \infty
\] 
then
\[
    \{X_n\} \text{ is tight}
\] 
}
Proof: $\forall K, \exists M $ such that $\phi(a) \ge K$ for $a \ge M$ this implies
\[
    \phi(a) \ge k 1_{a \ge M}
\] 
\[
    \mathbb{E}1_{|X_n| \ge M} \le \frac{1}{K} \mathbb{E} \phi(|X_n|) \le \frac{1}{K} \sup_n \mathbb{E} \phi(X_n) \le \epsilon
\] 
if we pick $K$ large enough\\
Example: 
\begin{enumerate}
    \item $X_n \sim B(n,\frac{1}{n})$. 
        \[
        \mathbb{E}|X_n| = \mathbb{E}X_n = 1
        \] 
        so the sequence is tight. The limit is Poisson(1)
    \item Let $\{\xi_k\}$ be i.i.d random variables with  $ \mathbb{E}\xi_i = 0$ $Var(\xi_i) =  1$
         \[
             X_{n} := \frac{\sum_{k=1}^{n}\xi_k}{\sqrt{n}}
        \] 
        \[
        \mathbb{E}X_n^2 = 1
        \] 
        so $\{X_n\}$ is tight. And converges to normal 0,1.
\end{enumerate}

\section{Uniformly integrable}
\[
    \lim_{M\to+\infty }(\sup_n \mathbb{E}X_n 1_{|X_n| \ge M}) = 0
\] 
Later we show that Uniform integrability and in probability if and only if in $L_1$

\section{Lemma} $\{X_n\}$. 
\emph{
    If $\exists \phi$ $ \mathbb{R} \rightarrow \mathbb{R}$ such that
    \[
    \lim_{a\to +\infty}\frac{\phi(a)}{a} = +\infty
    \] 
    and
    \[
        \sup_{n} \mathbb{E} \phi(|X_n|) < \infty
    \] 
    then $\{X_n\}$ is uniformly integrable
} 
Proof for all $k$ there exists $M$ such that $\phi(a) \ge ka$ for $a \ge M$ then $\phi(a) \ge ka 1_{a \ge m}$ and
\[
    \mathbb{E}|X_n|1_{|X_n|\ge M} = \frac{1}{k} \mathbb{E}k|X_n|1_{|X_n|\ge M} \le \frac{1}{k} \mathbb{E} \phi(|X_n|) \le \frac{1}{k} \sup_{n} \mathbb{E} \phi(|X_n|) \sim \epsilon
\] 

\section{Lemma UI + Prob implies $L_1$}
$L_1 \implies$ in probability we learned in 275 A. Let $X_n \rightarrow X$ in $L_1$ so all functions live in $ L_1$
\[
 \mathbb{E}|X_n - X| \rightarrow 0
\] 
so
for all  $\epsilon$ there exists $N$ such that $ \mathbb{E}|X_n-X| \le \epsilon$ for all $n \ge N$.\\
there exists $M$ such that $ \mathbb{E}|X|1_{|X| \ge M} \le \epsilon$ that implies that $ \mathbb{E}|X_n| 1_{|X_n| \ge M} \le 2\epsilon$ for all $n > N$.
Only finitely many dont satisfy this property so we can find a larger $M'$ that is satisfied for all of the  $X_n$
So the sequence is uniformly integrable.

\section{Friday}
UI + in Prob if and only if $L_1$.
$X_n \rightarrow X$ in $L_1$ iff $X_n \rightarrow X$ in prob and $X_n$ are U.I. Today we need to prove the forward direction.
First we want to prove that $X$ is in $L_1$
\[
    \forall \epsilon, \exists M \rightarrow \sup_n \mathbb{E}|X_n| 1_{|X_n \ge M|} \le \epsilon
\] 
but
\[
    \sup_n \mathbb{E}|X_n| 1_{|X_n \le M} \le M 
\] 
this implies
\[
\sup_n \mathbb{E}|X_n| < \infty
\] 
\[
    X_n \rightarrow X \text{ in prob} \implies \exists X_{n_k} \rightarrow X \text{ a.s}
\] 
Then by fatou's lemma
\[
    \mathbb{E}|X| \le \liminf_{k\to\infty } \mathbb{E}|X_{n_k}| < +\infty
\] 
so $X$ is integrable and $X \in L_1$. Now $X_n$ is U.I and  $X$ is in  $L_1$ implies
\[
    \forall \epsilon , \exists M \text{ s.t.} \quad \mathbb{E}|X_n| 1_{|X_n| \ge M} \le \epsilon \quad \mathbb{E}|X| 1_{|X| \ge M} \le \epsilon
\] 
goal
\[
\lim_{ } \mathbb{E}|X_n - X| = 0
\] 
\[
    \mathbb{E}|X_n - X| = \mathbb{E}|X_n - X|1_{|X_n|, |X| \le M} + \mathbb{E}|X_n - X|1_{|X_n| \text{ or } |X| \ge M}
\] 
\[
    Y_n = |X_n - X| 1_{|X_n|,|X| \le M}
\] 
\[
Y_n \rightarrow 0
\] 
in probability and $|Y_n| \le 2M$ almost sure. so
 \[
\lim_{n\to \infty} \mathbb{E}Y_n = 0
\] 
for the second part
\[
    \mathbb{E}|X_n-X| 1_{|X_n| \text{ or } |X| \ge M} \le \mathbb{E} 2 \max (|X_n|, |X|) 1_{\max(|X_n|,|X|) \ge M}
\] 
\[
    \le 2 \mathbb{E} |X_n| 1_{|X_n| \ge M} + 2 \mathbb{E} |X| 1_{|X| \ge M} \le 4 \epsilon
\] 
if we pick $M$ large enough. so we conclude that
\[
\mathbb{E}|X_n - X| \rightarrow 0
\] 

\section{Another Definition of UI}
\emph{
    \begin{enumerate}
        \item 
            \[
            \sup_n \mathbb{E}|X_n| < \infty
            \] 
        \item $\forall \epsilon$ $\exists \delta$ such that if $ \mathbb{P}(A) \le \delta$ then $\sup_n \mathbb{E} |X_n| 1_A \le \epsilon$
    \end{enumerate}
}

\section{Moment Method}
\emph{
    Lemma:  $\{X_n\}$. For all  $k \in \mathbb{N}$
    \[
    \mathbb{E}|X_n^{k}| \rightarrow \alpha_k < \infty
    \] 
    then for all $n_k$ there exists  $n_{k_m}$ and  $Y$ such that  $X_{n_{k_m}} \rightarrow^{d} Y$ and
    \[
    \mathbb{E}|Y^{k}| = \alpha_k
    \] 
    Furthermore, if there is only one random variable having $\{\alpha_k\}$ as its moments (uniqueness) then
    \[
        \exists Y \text{ such that } \mathbb{E}Y^{k} = \alpha_k \;\;\; X_n \rightarrow^{d} Y
    \] 
}
\[
G(t) = \sum_{k}^{}\frac{m_k}{k}e^{tk}
\] 
Proof Part A implies part B
 \[
     X_n \rightarrow^{d} X \iff \forall n_k \exists n_{k_m} \;\; X_{n_{k_m}} \rightarrow^{d} X
\] 
Proof part B $\{X_n\}$ are tight.  $ \sup_n \mathbb{E} X_n^2 <\infty$ so use $\phi(x) = x^2$ implies that $\{X_n\}$ is tight. which implies
the first part of the theorem of the existence of subsequential limits. We now need to prove that for all $k$  $  \mathbb{E}Y^{k} = \alpha_k$ 
We have 
\[
    Z_{n_{k_m}} \rightarrow Z \;\; \text{a.s} \qquad Z_{n_{k_m}} \sim X_{n_{k_m}}, Y \sim Z
\] 
\[
    \mathbb{E}Z_{n_{k_m}}^{k} \rightarrow \mathbb{E}z
\] 

\end{document}
